% This Source Code Form is subject to the terms of the Mozilla Public
% License, v. 2.0. If a copy of the MPL was not distributed with this
% file, You can obtain one at https://mozilla.org/MPL/2.0/.

\documentclass[a4paper,12pt]{article}

\usepackage[T2A]{fontenc}
\usepackage[utf8]{inputenc}
\usepackage[russian]{babel}
\usepackage{amsmath}
\usepackage{graphicx}
\usepackage{booktabs}
\usepackage{geometry}
\usepackage{tabularx}
\geometry{margin=2cm}

\begin{document}
	
	\subsubsection*{Опыт №1. Взаимодействие хлорида железа(III) и тиоцианата калия}
	
	При взаимодействии растворов \(\mathrm{FeCl_3}\) и \(\mathrm{KSCN}\) протекает обратимая реакция:
	\[
	\mathrm{Fe^{3+} + SCN^- \rightleftharpoons Fe(SCN)^{2+}}
	\]
	Образуется малодиссоциированное соединение --- тиоцианат железа(III), окрашивающее раствор в тёмно-красный цвет.  
	Степень окраски раствора служит показателем положения химического равновесия.
	
	Константа равновесия выражается как:
	\[
	K = \frac{[\mathrm{Fe(SCN)^{2+}}]}{[\mathrm{Fe^{3+}}][\mathrm{SCN^-}]}
	\]
	
	В 4 пробироки внесено по 6 капель разбавленных растворов \(\mathrm{FeCl_3}\) и \(\mathrm{KSCN}\) (с концентрацией \(C = 0{,}0025\) моль/л).  
	Растворы перемешаны и помещены в штатив. Одну пробирку оставили как \textbf{эталон}.  
	
	В остальные пробирки добавлено:
	\begin{enumerate}
		\item 1 капля насыщенного раствора \(\mathrm{FeCl_3}\);
		\item 1 капля насыщенного раствора \(\mathrm{KSCN}\);
		\item несколько кристалликов \(\mathrm{KCl}\).
	\end{enumerate}
	
	После добавления наблюдали изменение окраски по сравнению с эталоном.
	
	\subsubsection*{Наблюдения и выводы}
	
	\begin{center}
		\renewcommand{\arraystretch}{1.3} % немного увеличим высоту строк
		\setlength{\extrarowheight}{2pt}
		\begin{tabularx}{\textwidth}{|>{\raggedright\arraybackslash}X|
				>{\centering\arraybackslash}m{2.8cm}|
				>{\centering\arraybackslash}m{2.8cm}|
				>{\centering\arraybackslash}m{2.8cm}|
				>{\centering\arraybackslash}m{2.8cm}|}
			\hline
			\textbf{Изменения в равновесной системе} &
			\textbf{Добавление FeCl\textsubscript{3}} &
			\textbf{Добавление KSCN} &
			\textbf{Добавление KCl} \\ \hline
			Интенсивность окраски раствора &
			Увеличивается (тёмно-красный) &
			Увеличивается (тёмно-красный) &
			Ослабевает \\ \hline
			Смещение химического равновесия &
			Вправо &
			Вправо &
			Влево \\ \hline
		\end{tabularx}
	\end{center}
	
	\subsubsection*{Обсуждение результатов}
	
	\begin{itemize}
		\item Добавление \(\mathrm{FeCl_3}\) увеличивает концентрацию ионов \(\mathrm{Fe^{3+}}\), равновесие смещается вправо, усиливается окраска.
		\item Добавление \(\mathrm{KSCN}\) увеличивает концентрацию ионов \(\mathrm{SCN^-}\), равновесие также смещается вправо.
		\item Добавление \(\mathrm{KCl}\) повышает концентрацию ионов \(\mathrm{Cl^-}\), которые частично комплексируются с \(\mathrm{Fe^{3+}}\), снижая его активную концентрацию — равновесие смещается влево, раствор светлеет.
	\end{itemize}
	
	\subsubsection*{Опыт №2. Взаимодействие хлорида магния с гидроксидом аммония}
	
	При взаимодействии раствора хлорида магния и гидроксида аммония протекает обратимая реакция:
	\[
	\mathrm{MgCl_2 + 2NH_4OH \rightleftharpoons Mg(OH)_2 \downarrow + 2NH_4Cl}
	\]
	Гидроксид магния \(\mathrm{Mg(OH)_2}\) — малорастворимое соединение, которое выпадает в виде белого осадка.  
	Изменение концентрации ионов \(\mathrm{Mg^{2+}}\), \(\mathrm{OH^-}\) или \(\mathrm{NH_4^+}\) приводит к смещению равновесия в соответствии с принципом Ле Шателье.
	
	Выражение для константы равновесия:
	\[
	K = \frac{[\mathrm{NH_4^+}]^2 [\mathrm{Cl^-}]^2}{[\mathrm{Mg^{2+}}][\mathrm{OH^-}]^2}
	\]
	
	\subsubsection*{Наблюдения и выводы}
	
	\begin{center}
		\renewcommand{\arraystretch}{1.3}
		\setlength{\extrarowheight}{2pt}
		\begin{tabularx}{\textwidth}{|>{\raggedright\arraybackslash}X|
				>{\centering\arraybackslash}m{5cm}|}
			\hline
			\textbf{Добавление вещества} &
			\textbf{Наблюдаемое изменение (количество осадка Mg(OH)\textsubscript{2}) и направление смещения равновесия} \\ \hline
			Постепенное добавление \(\mathrm{NH_4OH}\) &
			Появление белого осадка \(\mathrm{Mg(OH)_2}\); равновесие смещается вправо. \\ \hline
			Добавление \(\mathrm{NH_4Cl}\) &
			Осадок частично растворяется; равновесие смещается влево (в сторону растворения). \\ \hline
			Повторное добавление \(\mathrm{NH_4OH}\) &
			Осадок вновь увеличивается; равновесие смещается вправо. \\ \hline
		\end{tabularx}
	\end{center}
	
	\subsubsection*{Обсуждение результатов}
	
	\begin{itemize}
		\item При добавлении гидроксида аммония увеличивается концентрация ионов \(\mathrm{OH^-}\), что способствует образованию осадка \(\mathrm{Mg(OH)_2}\); равновесие смещается вправо.
		\item При добавлении хлорида аммония увеличивается концентрация ионов \(\mathrm{NH_4^+}\), которые уменьшают концентрацию \(\mathrm{OH^-}\) за счёт протекания реакции нейтрализации; равновесие смещается влево, осадок частично растворяется.
		\item Повторное добавление гидроксида аммония снова повышает концентрацию \(\mathrm{OH^-}\), что возвращает равновесие вправо и осадок увеличивается.
	\end{itemize}

	\subsubsection*{Опыт №3. Взаимодействие диоксида азота и тетраоксида диазота}
	
	Изучить химическое равновесие между диоксидом азота и тетраоксидом диазота:
	\[
	\mathrm{N_2O_4 \rightleftharpoons 2NO_2}
	\]
	
	Диоксид азота (\(\mathrm{NO_2}\)) — бурый газ, а тетраоксид диазота (\(\mathrm{N_2O_4}\)) — бесцветный.  
	Интенсивность окраски газовой смеси зависит от соотношения этих соединений: при преобладании \(\mathrm{N_2O_4}\) смесь бледно-жёлтая, при преобладании \(\mathrm{NO_2}\) — бурого цвета.
	
	\textbf{Проведение опыта:}  
	Предварительно получают диоксид азота, действуя концентрированной азотной кислотой на медь:
	\[
	\mathrm{Cu + 4HNO_3 \rightarrow Cu(NO_3)_2 + 2NO_2 \uparrow + 2H_2O}
	\]
	Две пробирки с выделившимся бурым газом (\(\mathrm{NO_2}\)) соединяют резиновым шлангом, создавая общую систему. После установления равновесия между \(\mathrm{NO_2}\) и \(\mathrm{N_2O_4}\) одну пробирку помещают в стакан с холодной водой, а другую — в стакан с горячей водой.
	
	Через некоторое время наблюдают:
	\begin{itemize}
		\item В пробирке с горячей водой окраска \textbf{усиливается} (становится более бурой);
		\item В пробирке с холодной водой окраска \textbf{ослабевает} (становится бледно-жёлтой).
	\end{itemize}
	
	\subsubsection*{Наблюдения и выводы}
	
	\begin{center}
		\renewcommand{\arraystretch}{1.3}
		\setlength{\extrarowheight}{2pt}
		\begin{tabularx}{0.9\textwidth}{|>{\raggedright\arraybackslash}X|
				>{\centering\arraybackslash}m{4cm}|
				>{\centering\arraybackslash}m{4cm}|}
			\hline
			\textbf{Условия опыта} &
			\textbf{Пробирка в холодной воде} &
			\textbf{Пробирка в горячей воде} \\ \hline
			Наблюдаемое изменение &
			Раствор светлеет (окраска ослабевает) &
			Окраска усиливается (бурый цвет) \\ \hline
			Смещение равновесия &
			В сторону \(\mathrm{N_2O_4}\) &
			В сторону \(\mathrm{NO_2}\) \\ \hline
		\end{tabularx}
	\end{center}
	
	\subsubsection*{Обсуждение результатов}
	
	Реакция разложения \(\mathrm{N_2O_4}\) эндотермическая:
	\[
	\mathrm{N_2O_4 \rightleftharpoons 2NO_2 - Q}
	\]
	При повышении температуры (\(T \uparrow\)) равновесие смещается в сторону поглощения теплоты — то есть в сторону образования бурого \(\mathrm{NO_2}\).  
	При понижении температуры (\(T \downarrow\)) равновесие сдвигается влево, в сторону бесцветного \(\mathrm{N_2O_4}\).
	
	Таким образом, наблюдаемые изменения окраски подтверждают \textbf{принцип Ле Шателье}: при повышении температуры система стремится компенсировать воздействие, смещая равновесие в эндотермическом направлении.
	
	\begin{center}
		\[
		K = \frac{P^2_{\mathrm{NO_2}}}{P_{\mathrm{N_2O_4}}}
		= \frac{\varphi^2_{\mathrm{NO_2}}}{\varphi_{\mathrm{N_2O_4}}}
		\]
	\end{center}
	
\end{document}
